% ======================================================================
\section{Light-cone wave functions and the coefficients of \texorpdfstring{$\Omega$}{Omega}}\label{app:wavefunctions}
% ======================================================================

The coefficients of $\Omega$ listed in Section~\ref{sec:coefficients} are obtained by computing the action of $\Omega$ on the soft vacuum and on a one-gluon state in two ways: from the expansion~\eqref{eq:OmegaExp}, and from LCPT, Eq.~\eqref{eq:lcpt}, with the matrix elements of Appendix~\ref{app:matrixelements}. The two are equated sector by sector in Fock space. In this appendix we give the Fock components, the diagrams from which they arise, and the steps that lead from the components to the coefficients. In the diagrams the horizontal line denotes the valence charges, curly lines soft gluons, and vertical dashed lines the intermediate states whose energies enter the denominators; a vertical bar on a line denotes an instantaneous interaction. Momenta are written in momentum space; the coordinate forms follow with the Fourier transforms of Appendix~\ref{app:fourier}.

\subsection{Matching conditions}

From Eq.~\eqref{eq:OmegaExp} and the normalisation of the states,
\begin{align}\label{eq:Omegavac}
\Omega\ket0=\mf N\ket0+(2\pi)^{3/2}\int\frac{d^3p}{(2\pi)^3}\Big[\mf A^{(3)b}_{2j}(p)-\mf A^{(1)\dagger b}_{1j}(p)\Big]\ket{g^b_j(p)}+(2\pi)^3\int\frac{d^3p}{(2\pi)^3}\frac{d^3q}{(2\pi)^3}\,\mf B^{cb}_{2kj}(p,q)\ket{g^b_j(p)\,g^c_k(q)}+\dots
\end{align}
and
\begin{align}\label{eq:Omegag}
\Omega\ket{g^a_i(k)}=\ &\frac{1}{(2\pi)^{3/2}}\Big[\mf A^{(1)a}_{1i}(k)+\mf A^{(3)a}_{1i}(k)\Big]\ket0+\mf N\ket{g^a_i(k)}+\int\frac{d^3p}{(2\pi)^3}\,\mf B^{ba}_{3ji}(k,p)\ket{g^b_j(p)}
\notag\\&
+(2\pi)^{3/2}\int\frac{d^3p}{(2\pi)^3}\Big[\mf A^{(3)b}_{2j}(p)-\mf A^{(1)\dagger b}_{1j}(p)\Big]\ket{g^a_i(k)\,g^b_j(p)}
+(2\pi)^{3/2}\int\frac{d^3p}{(2\pi)^3}\,\mf C^{bca}_{1jki}(k,p)\ket{g^b_j(p)\,g^c_k(k-p)}+\dots,
\end{align}
where the dots denote states with more gluons. The vacuum sector fixes $\mf N$, $\mf A^{(1)}$, $\mf A^{(3)}_2$ and $\mf B_2$; the one-gluon sector fixes $\mf B_3$ and $\mf C_1$, and reproduces $\mf A^{(1)}_1$ and the product $\mf A\,$(one-gluon state) as a consistency check.

\subsection{The soft vacuum at order \texorpdfstring{$g$}{g} and the normalisation}

At order $g$ the valence charges emit one soft gluon (Figure~\ref{fig:vac_g}):
\begin{align}\label{eq:PsiLO}
\ket{\Psi^{\rm LO}}=-\int_\Lambda^\vee dk^+\!\int d^2k\,\ket{g^a_i(k)}\frac{\braket{g^a_i(k)|H_g|0}}{E_g(k)}
=-\int_\Lambda^\vee dk^+\!\int d^2k\,\frac{g\,k^i\rho^a(-\v k)}{2\pi^{3/2}\sqrt{k^+}\,\v k^2}\ket{g^a_i(k)}
=\frac{ig}{2\pi^{3/2}}\int_\Lambda^\vee\frac{dk^+}{\sqrt{k^+}}\int_{x,z}\Kv^i(\v x-\v z)\rho^a(\v x)\ket{g^a_i(k^+,\v z)}\,.
\end{align}
Comparison with Eq.~\eqref{eq:Omegavac} gives $-\mf A^{(1)a}_{2i}(k^+,\v k)=\mf A^{(1)\dagger a}_{1i}(k^+,\v k)=\sqrt2\,g\,k^i\rho^a(-\v k)/(\sqrt{k^+}\v k^2)$, i.e.\ Eq.~\eqref{eq:A1}. The norm of the state to order $g^2$ fixes
\begin{align}
\mf N=1-\frac12\braket{\Psi^{\rm LO}|\Psi^{\rm LO}}=1-\frac{g^2}{8\pi^3}\int_\Lambda^\vee\frac{dp^+}{p^+}\int d^2p\,\frac{\rho^a(-\v p)\rho^a(\v p)}{\v p^2}\,,
\end{align}
which is Eq.~\eqref{eq:N2}.

\begin{figure}[t]
\centering
\tikzfig{vacuumtog}
\caption{Emission of a soft gluon by the valence charges at order $g$.}
\label{fig:vac_g}
\end{figure}

\subsection{Two soft gluons at order \texorpdfstring{$g^2$}{g2}}

Three diagrams contribute (Figure~\ref{fig:vac_gg}): the emission of two gluons by the valence charges, with the matrix elements~\eqref{eq:m1} and~\eqref{eq:m2},
\begin{align}
\ket{\Psi^{1}_{gg}}=\frac12\int dq^+dk^+dp^+\!\int d^2q\,d^2k\,d^2p\,\ket{g^b_j(k-p)\,g^c_l(p)}\frac{\braket{g^b_j(k-p)\,g^c_l(p)|H_g|g^a_i(q)}\braket{g^a_i(q)|H_g|0}}{E_{gg}(p,k-p)\,E_g(q)}\,,
\end{align}
the emission of one gluon which splits into two through the triple-gluon vertex~\eqref{eq:m3},
\begin{align}
\ket{\Psi^{2}_{gg}}=\frac12\int dk^+dp^+dq^+\!\int d^2k\,d^2p\,d^2q\,\ket{g^b_j(q)\,g^c_l(p)}\frac{\braket{g^b_j(q)\,g^c_l(p)|H_{ggg}|g^a_i(k)}\braket{g^a_i(k)|H_g|0}}{E_{gg}(p,q)\,E_g(k)}\,,
\end{align}
and the instantaneous emission of the pair~\eqref{eq:m4},
\begin{align}
\ket{\Psi^{3}_{gg}}=-\frac12\int dk^+dp^+\!\int d^2k\,d^2p\,\ket{g^b_i(k-p)\,g^c_j(p)}\frac{\braket{g^b_i(k-p)\,g^c_j(p)|H_{gg\text{-inst}}|0}}{E_{gg}(k-p,p)}\,.
\end{align}
In the first diagram the product of the two charges is decomposed as
\begin{align}\label{eq:decomp}
\rho^b(-\v k+\v p)\,\rho^a(-\v p)=\frac{i}{2}f^{bad}\rho^d(-\v k)+\frac12\big\{\rho^b(-\v k+\v p),\rho^a(-\v p)\big\}\,.
\end{align}
The commutator term is linear in $\rho$ and is combined with the other two diagrams. The anticommutator term describes two independent emissions. Comparing with Eq.~\eqref{eq:Omegavac} gives the coefficient
\begin{align}\label{eq:B2mom}
    \mf{B}^{ba}_{2ji}(k^{+},k;p^{+},p) & =  \frac{ig^{2}f^{abc}\,p^{j}k^{i}\rho^{c}(-p-k)\, \sqrt{p^{+}k^{+}}}{2\,\left(p^{+}k^{2}+k^{+}p^{2}\right)}\,\left[\frac{1}{k^{+}\,p^{2}} - \frac{1}{k^{2}\,p^{+}}\right]\,\Theta\left(\vee-k^{+}-p^{+}\right)\notag \\ & + \frac{g^{2}p^{j}k^{i}\,\{\rho^{b}(-p),\rho^{a}(-k)\}}{2\,p^{2}k^{2}\, \sqrt{p^{+}k^{+}}}\,\Theta\left(\vee-k^{+}\right)\Theta\left(\vee-p^{+}\right)\notag \\ & -\,\frac{ig^{2}f^{abc}\rho^{c}(-p-k)\,(p+k)^{l}\,\left(p^{m}k^{+}-p^{+}k^{m}\right)}{(p+k)^{2}\,\sqrt{p^{+}k^{+}}\,(p^{+}+k^{+})\,\left(\frac{p^{2}}{p^{+}}+\frac{k^{2}}{k^{+}}\right)}\,\left[\frac{\delta_{lm}\delta_{ji}}{(p^{+}+k^{+})} - \frac{\delta_{im}\delta_{jl}}{k^{+}} - \frac{\delta_{jm}\delta_{il}}{p^{+}}\right]\notag \\ & \qquad\times\Theta\left(\vee-k^{+}-p^{+}\right)\notag \\ & -\,\frac{ig^{2}f^{abc}\rho^{c}(-p-k)\,(p^{+}-k^{+})\delta_{ji}}{2\,\sqrt{p^{+}k^{+}}\,(p^{+}+k^{+})^{2}\,\left(\frac{k^{2}}{k^{+}}+\frac{p^{2}}{p^{+}}\right)}\,\Theta\left(\vee-k^{+}-p^{+}\right)
\end{align} [$\Theta$ replaced; the support follows from the kinematics ($k^+,p^+>\Lambda$ throughout). Third term (triple vertex): the parent gluon carries $k^++p^+$ and is itself a soft mode, hence $\Theta(\vee-k^+-p^+)$; the instantaneous fourth term has the same support in Eq.~\eqref{ggthree} (a convention; no particle carries $k^++p^+$ there). Second term ($\{\rho,\rho\}$): the two gluons are emitted independently by the valence charges, so $\Lambda<k^+<\vee$ and $\Lambda<p^+<\vee$ with no condition on the sum, as for the same intermediate state in Eq.~\eqref{3rho}; with $k^++p^+<\vee$ the four-$\rho$ terms of the cross section would not cancel for commuting $\rho$ (coherent state), see the note after Eq.~\eqref{eq:CS1}. The first term is the commutator part of the same diagram, so by the same kinematics its support is also the full square; it is left unchanged here for the author to decide. $\Theta(k^+-\Lambda-p^+)$ is not symmetric under exchange of the two gluons and removes $p^+>k^+$.]\par\noindent 
where the first, third and fourth terms come from the commutator part of the first diagram, from the triple-gluon vertex and from the instantaneous emission, respectively, and the second term from the anticommutator part. The supports are discussed in Section~\ref{sec:coefficients}. With the Fourier transforms~\eqref{eq:FT_i2} and~\eqref{eq:FT_i7} this becomes Eq.~\eqref{eq:B2}.

\begin{figure}[t]
\centering
\tikzfig{vacuumto2g2rho}\qquad\tikzfig{vacuumto2g}\qquad\tikzfig{vacuumto2ginst}
\caption{Two soft gluons at order $g^2$: two emissions by the valence charges, emission followed by splitting, and instantaneous emission.}
\label{fig:vac_gg}
\end{figure}

\subsection{One soft gluon at order \texorpdfstring{$g^3$}{g3}}

\paragraph{Diagrams.} The one-gluon component at order $g^3$ receives contributions linear, quadratic and cubic in the charges.
\begin{itemize}
\item The self-energy of the emitted gluon through two triple-gluon vertices (Figure~\ref{fig:vac_g3_onerho}),
\begin{align}
\ket{\Psi^{1}_{g\rho}}=-\frac12\int dk^+dp^+dq^+dr^+\!\int d^2k\,d^2p\,d^2q\,d^2r\,\ket{g^d_l(r)}\frac{\braket{g^d_l(r)|H_{ggg}|g^c_k(q)g^b_j(p)}\braket{g^c_k(q)g^b_j(p)|H_{ggg}|g^a_i(k)}\braket{g^a_i(k)|H_g|0}}{E_g(r)\,E_{gg}(p,q)\,E_g(k)}\,.
\end{align}
It is linear in $\rho$ and is the only diagram with an ultraviolet divergence of the transverse loop integral.
\item The diagrams with two valence vertices (Figure~\ref{fig:vac_g3_tworho}): the emitted gluon interacts instantaneously with the valence charges, gaining or losing longitudinal momentum; two gluons emitted by the valence charges recombine through the triple-gluon vertex; two gluons are emitted, one instantaneously, and one of them is absorbed by the valence charges; a gluon splits into two, one of which is absorbed by the valence charges. With the decomposition~\eqref{eq:decomp} each of them gives a term linear in $\rho$ (commutator) and a term $\propto f^{abc}\{\rho^b,\rho^c\}$.
\item The diagrams with three valence vertices (Figure~\ref{fig:vac_g3_threerho}): a gluon is emitted and reabsorbed by the valence charges before or after the emission of the observed gluon, in the three time orderings, and the corresponding normalisation term of Eq.~\eqref{eq:lcpt}.
\end{itemize}

\begin{figure}[t]
\centering
\tikzfig{loop}
\caption{Self-energy of the emitted gluon at order $g^3$.}
\label{fig:vac_g3_onerho}
\end{figure}

\begin{figure}[t]
\centering
\tikzfig{Misha_2rho_single_g}\qquad\tikzfig{Misha_Out_2rho_single_g}\qquad\tikzfig{Misha_2rho_3_single_g}\\[4mm]
\tikzfig{Misha_2rho_insta_single_g}\qquad\tikzfig{Misha_2rho_loop_single_g}
\caption{Diagrams with two valence vertices contributing to the one-gluon component at order $g^3$.}
\label{fig:vac_g3_tworho}
\end{figure}

\begin{figure}[t]
\centering
\tikzfig{Misha_2rho_loopVal1_single_g}\qquad\tikzfig{Misha_2rho_loopVal2_single_g}\qquad\tikzfig{Misha_2rho_loopVal3_single_g}
\caption{Diagrams with three valence vertices contributing to the one-gluon component at order $g^3$.}
\label{fig:vac_g3_threerho}
\end{figure}

\paragraph{Self-energy.} After the transverse loop integral in $d=2-\epsilon$ dimensions, with $p^+=\xi k^+$,
\begin{align}
\ket{\Psi^{1}_{g\rho}}=-\int_\Lambda^\vee\frac{dk^+}{\sqrt{k^+}}\int d^2k\int_{\Lambda/k^+}^{1-\Lambda/k^+}d\xi\,\frac{g^3N_c\,\rho^a(-\v k)\,k^i}{16\pi^{7/2}\,\v k^2}\Big[-2+\xi(1-\xi)+\frac{1}{\xi(1-\xi)}\Big]\Big[-\frac2\epsilon+\ln\frac{\v k^2\xi(1-\xi)}{\mu^2_{\MS}}\Big]\ket{g^a_i(k^+,\v k)}\,.
\end{align}
With $\delta=\Lambda/k^+\to0$,
\begin{align}\label{eq:xiintegrals}
\int_\delta^{1-\delta}d\xi\Big[-2+\xi(1-\xi)+\frac{1}{\xi(1-\xi)}\Big]=-\frac{11}{6}-2\ln\delta\,,\qquad
\int_0^1d\xi\,\big[-2+\xi(1-\xi)\big]\ln\big(\xi(1-\xi)\big)=\frac{67}{18}\,,\qquad
\int_\delta^{1-\delta}d\xi\,\frac{\ln(\xi(1-\xi))}{\xi(1-\xi)}=-\ln^2\delta-\frac{\pi^2}{3}\,,
\end{align}
so that
\begin{align}
\ket{\Psi^{1}_{g\rho}}=\int_\Lambda^\vee\frac{dk^+}{\sqrt{k^+}}\int d^2k\,\frac{g^3N_c\,\rho^a(-\v k)\,k^i}{32\pi^{7/2}\,\v k^2}\Big[\Big(\frac{11}{3}+4\ln\frac{\Lambda}{k^+}\Big)\Big(-\frac2\epsilon+\ln\frac{\v k^2}{\mu^2_{\MS}}\Big)+2\ln^2\frac{\Lambda}{k^+}-\frac{67}{9}+\frac{2\pi^2}{3}\Big]\ket{g^a_i(k^+,\v k)}\,.
\end{align}
This is the origin of the coefficient $\frac{11}{3}=\beta_0/N_c$ and of the constant $\frac{67}{9}$.

\paragraph{Sum of the one-$\rho$ terms.} Adding the commutator parts of the diagrams with two and three valence vertices, the one-gluon component linear in $\rho$ is
\begin{align}\label{eq:Psigrho}
\ket{\Psi_{g\rho}}=\ &\int_\Lambda^\vee\frac{dk^+}{\sqrt{k^+}}\int d^2k\,\frac{g^3N_c\,\rho^a(-\v k)\,k^i}{32\pi^{7/2}\,\v k^2}\Big[\Big(-\frac2\epsilon+\ln\frac{\v k^2}{\mu^2_{\MS}}\Big)\Big(\frac{11}{3}+\ln\frac{\Lambda}{k^+}-2\ln\frac{\vee}{\Lambda}\Big)+\frac{\pi^2}{3}-\frac{67}{9}+\ln^2\frac{\Lambda}{k^+}-\ln^2\frac{\vee}{k^+}\Big]\ket{g^a_i(k^+,\v k)}
\notag\\&
+\int_\Lambda^\vee\frac{dk^+}{\sqrt{k^+}}\int d^2k\int_{\Lambda/k^+}^{1-\Lambda/k^+}d\xi\,\frac{g^3N_c\,\rho^a(-\v k)\,k^i}{32\pi^{7/2}\sqrt\xi(1-\xi)\,\v k^2}\Big[(\xi-2)\ln\xi+(\xi+1)\ln(1-\xi)-(\xi-2)\Big(-\frac2\epsilon+\ln\frac{\v k^2}{\mu^2_{\MS}}\Big)\Big]\ket{g^a_i(\xi k^+,\v k)}\,.
\end{align}
The second line, in which the observed gluon carries the momentum fraction $\xi$ of the parent, comes from the diagrams in which a gluon of momentum $k^+$ transfers part of its longitudinal momentum to the valence charges. To read off the coefficient of a gluon of momentum $k^+$ we relabel $\xi k^+\to k^+$ in this term; the parent momentum is then $k^+/\xi$, and the conditions $\Lambda<k^+/\xi<\vee$, $\xi<1-\Lambda\xi/k^+$ give $k^+/\vee<\xi<k^+/(k^++\Lambda)$. The $\xi$ integrals,
\begin{align}
\int_{k^+/\vee}^{k^+/(k^++\Lambda)}d\xi\,\frac{2-\xi}{\xi(1-\xi)}&=2\ln\frac{\vee}{k^+}+\ln\frac{k^+}{\Lambda}+\ln\Big(1-\frac{k^+}{\vee}\Big)+O\Big(\frac{\Lambda}{k^+}\Big)\,,\notag\\
\int_{k^+/\vee}^{k^+/(k^++\Lambda)}d\xi\,\frac{(\xi-2)\ln\xi+(\xi+1)\ln(1-\xi)}{\xi(1-\xi)}&=-\ln^2\frac{\Lambda}{k^+}+\ln^2\frac{\vee}{k^+}+\ln^2\Big(1-\frac{k^+}{\vee}\Big)+\ln\Big(1-\frac{k^+}{\vee}\Big)\ln\frac{\vee}{k^+}+O\Big(\frac{\Lambda}{k^+}\ln\frac{\Lambda}{k^+}\Big)\,,
\end{align}
add $\ln(\vee/k^+-1)+\ln(\vee/\Lambda)$ to the coefficient of the pole and convert the constant of the first line into $\frac{\pi^2}{3}-\frac{67}{9}+\ln^2(1-k^+/\vee)+\ln(1-k^+/\vee)\ln(\vee/k^+)$. The double logarithms of the cut-offs cancel, and the result is Eq.~\eqref{eq:A3onerho}.

\paragraph{Two-$\rho$ terms.} The anticommutator parts of the diagrams with two and three valence vertices give
\begin{align}\label{eq:Psigrhorho}
\ket{\Psi_{g\rho\rho}}=\int_\Lambda^\vee\frac{dk^+}{\sqrt{k^+}}\int d^2k\,d^2p\,\frac{ig^3f^{abc}\{\rho^b(-\v p),\rho^c(-\v k+\v p)\}}{32\pi^{9/2}\,\v p^2\,\v k^2}\Bigg[&\ln\frac{\Lambda}{k^+}\,\frac{\v k^2p^i-(\v k\cdot\v p)k^i}{(\v k-\v p)^2}-2k^i\ln\frac{\vee}{\Lambda}-2k^i\ln\frac{\Lambda\v k^2+k^+\v p^2}{\vee\v k^2+k^+\v p^2}
\notag\\&
+\frac{(2\v p^2-2\v k\cdot\v p+\v k^2)k^i+\v k^2p^i}{(\v k-\v p)^2}\ln\Big[1+\Big(\frac{\vee}{k^+}-1\Big)\frac{\v k^2}{(\v k-\v p)^2}\Big]
\notag\\&
-\frac{\v p^2k^i+\v k^2p^i}{(\v k-\v p)^2}\ln\frac{\vee-k^+}{\Lambda}-k^i\ln\frac{\vee}{k^+}\Bigg]\ket{g^a_i(k)}\,,
\end{align}
after the $\xi$ integration. The terms proportional to $k^+/\Lambda$ of the individual instantaneous diagrams cancel in the sum. Setting $\Lambda\to0$ where possible and transforming to coordinate space gives Eq.~\eqref{eq:A3tworho}.

\paragraph{Three-$\rho$ terms.} The diagrams with three valence vertices give, in addition to commutator terms,
\begin{align}
\ket{\Psi^{6c}_{g\rho\rho\rho}}&=\int_\Lambda^\vee dk^+\!\int d^2k\,d^2p\,\frac{g^3\,k^i\,\rho^b(\v p)\rho^b(-\v p)\rho^a(-\v k)\sqrt{k^+}}{8\pi^{9/2}\,\v k^4}\Big[\frac1\vee-\frac1\Lambda\Big]\ket{g^a_i(k^+,\v k)}\,,\notag\\
\ket{\Psi^{7}_{g\rho\rho\rho}}&=\int_\Lambda^\vee\frac{dk^+}{\sqrt{k^+}}\int d^2k\,d^2p\,\frac{g^3\,k^i\,\rho^a(-\v k)\rho^b(\v p)\rho^b(-\v p)}{16\pi^{9/2}\,\v k^2}\Big[\frac{1}{\v p^2}\ln\frac{\vee}{\Lambda}+\frac{2k^+}{\v k^2}\Big(\frac1\Lambda-\frac1\vee\Big)\Big]\ket{g^a_i(k^+,\v k)}\,.
\end{align}
The second line is the normalisation term of Eq.~\eqref{eq:lcpt}, with the single-emission matrix element to the left of the norm factor. The terms proportional to $1/\Lambda$ and $1/\vee$ cancel exactly between the two lines, although the charges stand in different orders: $\int d^2p\,\rho^b(\v p)\rho^b(-\v p)$ without the weight $1/\v p^2$ is local in coordinate space and commutes with $\rho^a$. What remains is
\begin{align}\label{eq:Psigrhorhorho}
\ket{\Psi_{g\rho\rho\rho}}=\ln\frac{\vee}{\Lambda}\int_\Lambda^\vee\frac{dk^+}{\sqrt{k^+}}\int d^2k\,d^2p\,\frac{g^3\,\rho^a(-\v k)\,\rho^b(\v p)\,\rho^b(-\v p)\,k^i}{16\pi^{9/2}\,\v k^2\,\v p^2}\ket{g^a_i(k^+,\v k)}\,,
\end{align}
which gives $[\mf A^{(3)}_2]_{\rho\rho\rho}=\mf A^{(1)}_2\mf N_2$, Eq.~\eqref{eq:A3threerho}.

\subsection{One soft gluon in the initial state}

The coefficients $\mf B_3$ and $\mf C_1$ are fixed by the dressing of a one-gluon state $\ket{g^a_i(k)}$ at order $g^2$ and $g$, respectively. The diagrams for the transition into a one-gluon state are shown in Figure~\ref{fig:g_g}: the instantaneous interaction with the valence charges, for $p^+>k^+$ and $p^+<k^+$; the emission of a second gluon by the valence charges and the absorption of the first by the triple-gluon vertex, or vice versa; and the diagrams with two valence vertices, in which the first gluon is absorbed and a second one emitted, or the second one emitted before the first is absorbed. The splitting of the gluon into two soft gluons is shown in Figure~\ref{fig:g_gg}. The energy denominators are $E_g(k)-E_n$ (Appendix~\ref{app:lcptrules}). Comparison with Eq.~\eqref{eq:Omegag} gives
\begin{align}\label{eq:B3mom}
    \mf{B}^{ba}_{3ji}(k^{+},k; p^{+},p) & =  \frac{ig^2 f^{abc}\,\sqrt{p^{+}k^{+}} (k^++p^+) } {(p^+-k^+)^2} \frac{\rho^c(-\v{p}+\v{k})}{p^{+}\v{k}^2 - k^{+}\v{p}^2} \delta_{ij}\left[\Theta\left(p^{+} - k^{+} - \Lambda\right)+\Theta\left(k^{+}-\Lambda-p^{+}\right)\right]\notag \\ & -\frac{g^{2}\,k^{i}p^{j}}{k^{2}p^{2}\,\sqrt{k^{+}p^{+}}}\left[\{\rho^{a}(k),\rho^{b}(-p)\}-if^{abc}\rho^{c}(k-p)\frac{p^{2}k^{+}+k^{2}p^{+}}{p^{2}k^{+}-k^{2}p^{+}}\right]\notag \\ & +\frac{2i\,g^{2}f^{abc}\,\rho^{c}(k-p)(k-p)^{l}\,\sqrt{k^{+}p^{+}}}{(k^{+}-p^{+})^{2}}\times \frac{\left(k^{m}p^{+}-k^{+}p^{m}\right)}{\left(k^{2}p^{+}-p^{2}k^{+}\right)}\,\left[\frac{\delta_{jm}\delta_{il}}{p^{+}}-\frac{\delta_{im}\delta_{jl}}{k^{+}}-\frac{\delta_{lm}\delta_{ij}}{(p^{+}-k^{+})}\right]\times \notag \\ & \left[\frac{(p^{+}-k^{+})}{(k-p)^{2}}\,\Theta(p^{+}-k^{+}-\Lambda)+\frac{1}{\left(\frac{k^{2}}{k^{+}}-\frac{(k-p)^{2}}{(k^{+}-p^{+})}-\frac{p^{2}}{p^{+}}\right)}\,\Theta(k^{+}-\Lambda-p^{+})\right]
\end{align} [Sign of the non-instantaneous $p^+>k^+$ term fixed (Eqs.~\eqref{onegrho12} and \eqref{Eq:onegrhoX1} carry the wrong overall sign): it follows from \eqref{m2}, \eqref{m3}, from the cancellation of the $1/(1-\xi)^2$ singularity, and from the $a^\dagger a$ unitarity constraint.]\par\noindent 
and
\begin{align}\label{eq:C1mom}
    \mf{C}^{bca}_{1jki}(k^{+},k; p^{+},p) & = -\,\frac{igf^{abc}\, \left(p^{m}k^{+}-p^{+}k^{m}\right)\,\sqrt{k^{+}p^{+}(k^{+}-p^{+})}}{\sqrt{2}\,\left(pk^{+}-p^{+}k\right)^{2}}\,\left[\frac{\delta_{im}\delta_{jk}}{k^{+}}-\frac{\delta_{ik}\delta_{mj}}{p^{+}}-\frac{\delta_{km}\delta_{ij}}{(k^{+}-p^{+})}\right]\,\Theta\left(k^{+}-\Lambda-p^{+}\right) 
\end{align}
The double poles $1/(k^+-p^+)^2$ of the instantaneous term and of the term with the triple-gluon vertex cancel each other at $|p^+-k^+|\to\Lambda$, so that no term proportional to $k^+/\Lambda$ arises. The coordinate forms, obtained with the Fourier transforms~\eqref{eq:FT_i3}--\eqref{eq:FT_i5}, are Eqs.~\eqref{eq:B3} and~\eqref{eq:C1}.

\begin{figure}[t]
\centering
\tikzfig{gtog_inst}\qquad\tikzfig{gandgtog}\qquad\tikzfig{gto3g_torho}\\[4mm]
\tikzfig{gtovactog_tworho}\qquad\tikzfig{gg_tworho}
\caption{Transition of a soft gluon into a soft gluon through the interaction with the valence charges at order $g^2$.}
\label{fig:g_g}
\end{figure}

\begin{figure}[t]
\centering
\tikzfig{g_togg}\qquad\tikzfig{gtog_plusg}
\caption{Splitting of a soft gluon into two soft gluons ($\rho$-independent, order $g$), and emission of a second soft gluon by the valence charges in the presence of the first one.}
\label{fig:g_gg}
\end{figure}
